\section{Conclusion}\label{sec:conclusion}

For long-tailed small rice pests run through standard overlapping sliced inference, the slice boundary is a secondary contributor to the misses it is usually blamed for. A pre-registered, group-disjoint, size-controlled audit shows that literal truncation is negligible and that the audited tail-small objects are lost predominantly at candidate formation, before any merge, threshold or cap can act, whereas same-size head-class pests are lost after a candidate exists and are therefore within reach of a better read-out. Because the majority of the loss sits where no post-processing can reach it, same-evidence rescoring recovers none of the tail-small misses and a class-conditional read-out adds nothing over class-by-size calibration; the admissible fix for these tail-small misses lies in the detector and its training signal for rare small classes rather than in the read-out. The audit reframes the problem rather than solving it, and its value is that it prescribes the class of fix before one is built.

The most important limitation is the size of the confirmatory endpoint, 38 tail-small objects dominated by one leafhopper species, which makes the primary result a pooled proportion rather than a per-class law and leaves its seed-to-seed stability unmeasured under the single-seed budget. The concrete next step is to repeat the audit with several training seeds and on a native-resolution, naturally long-tailed rice-pest source once a license-clear one is available, which would turn the pooled finding into per-class statements and test directly whether the candidate-formation bottleneck persists when the small objects are not the product of downscaling.
